\section{Zusammenfassung und Fazit}

Aufgrund des Vorhandenseins einer Filament-Rolle konnte zeitweise nur ein 3D-Drucker für die Fertigung genutzt werden, wodurch sich die benötigte Druckzeit erheblich verlängerte. Um ein einheitliches Erscheinungsbild zu erreichen, wurden die Bauteile, mit Ausnahme der Schrift, mit demselben Filament gefertigt. Dadurch lagen die Farben der Halterungen und der darin befestigten Bauteile teilweise sehr nah beieinander. Dies erschwert beispielsweise beim Sicherungsring die optische Unterscheidung zwischen Bauteil und Halterung (siehe Abb. \ref{fig:Halterung mit Sicherungsring}).

\begin{figure}[H]
	\centering
	\includegraphics[width=0.4\linewidth]{Halterung mit Sicherungsring.jpeg}
	\caption{Halterung mit Sicherungsring}
	\label{fig:Halterung mit Sicherungsring}
\end{figure}

Auch die Platten mit Beschriftungen und QR-Codes bieten noch Verbesserungspotenzial. Die Qualität der gedruckten Schrift war nicht zufriedenstellend. Zudem konnte der QR-Code nicht direkt auf die Platten gedruckt werden und musste daher zusätzlich ausgedruckt und angebracht werden.

Neben diesen Einschränkungen konnte die Aufgabe insgesamt zufriedenstellend erfüllt werden. Es wurden modulare Halterungen erstellt, die bei einem Wechsel der Ausstellung einfach angepasst werden können. Dadurch ist es möglich, einzelne Bauteile auszutauschen oder die Halterungen anzupassen, ohne die Konstruktion vollständig neu erstellen zu müssen.

